% The method figure. Panel (a) is why the comparison cannot be made on the span;
% panel (b) is what the ratio is made of.
%
% Layout note: every row is built right-to-left from the merge line, so the rows
% are contiguous and the shared suffix starts at the same x in all four. An
% earlier version left-aligned the spans, which opened a gap before the suffix
% and read as missing tokens -- the opposite of the panel's point, which is that
% the views differ in span *width* and agree from the merge position onward.
%
% The two image strips are the actual Tier B stimulus at the frozen render
% geometry, written by scripts/make_figure_assets.py.
\begin{figure}[t]
\centering
\begin{tikzpicture}[
    font=\small,
    tok/.style={draw=black!45, rounded corners=1pt, minimum height=4.4mm,
                inner xsep=2.5pt, fill=black!4, font=\scriptsize\ttfamily},
    span/.style={tok, fill=orange!16, draw=orange!65},
    vtok/.style={draw=blue!55, fill=blue!14, rounded corners=0.4pt,
                 minimum height=1.5mm, minimum width=2.55mm, inner sep=0pt},
    lbl/.style={font=\footnotesize, anchor=east},
    >=stealth,
]
\def\imgw{1.7cm}

% ---------------------------------------------------------------- panel (a) ---
\node[font=\footnotesize\bfseries, anchor=west] at (-3.75, 1.02) {(a)};

% --- text views: prefix and span hang leftward from the merge line at x = 0 ---
\foreach \y/\name/\word in {0/T/wisdom, -0.78/Tc/insight} {
  \node[span, anchor=east] (s\name) at (0, \y) {\word};
  \node[tok, anchor=east] (p\name) at ($(s\name.west)+(-1pt,0)$) {She\ wrote};
}
\node[lbl] at ($(pT.west)+(-4pt,0)$) {$V_\mathrm{T}$};
\node[lbl] at ($(pTc.west)+(-4pt,0)$) {$V_\mathrm{T}'$};

% --- image views: the rendered strip, with the six visual tokens beneath it ---
\foreach \y/\name/\file in {-1.85/I/span_primary, -2.95/Ic/span_control} {
  \node[anchor=east, inner sep=0pt, draw=blue!45, line width=0.3pt]
       (i\name) at (0, \y+0.12) {\includegraphics[width=\imgw]{figures/\file.png}};
  \foreach \i in {0,...,5} {
    \node[vtok, anchor=east] at ($(i\name.south east)+(-\i*2.65mm, -1.6mm)$) {};
  }
  \node[tok, anchor=east] (p\name) at ($(i\name.west)+(-1pt,0)$) {She\ wrote};
}
\node[lbl] at ($(pI.west)+(-4pt,0)$) {$V_\mathrm{I}$};
\node[lbl] at ($(pIc.west)+(-4pt,0)$) {$V_\mathrm{I}'$};

% --- the shared, terminal suffix: identical and identically placed ----------
\foreach \y in {0, -0.78, -1.85, -2.95} {
  \node[tok, anchor=west, fill=black!9] at (2pt, \y) {on\ the\ page.};
}

% --- the merge position ----------------------------------------------------
\draw[dashed, red!70, thick] (0, 0.62) -- (0, -3.34);
\node[font=\footnotesize, red!75, anchor=south] at (0.92, 0.66) {merge position $m$};

\draw[decorate, decoration={brace, amplitude=3pt, mirror}, black!55]
     (-1.75, -3.40) -- (-0.03, -3.40);
\node[font=\scriptsize, anchor=north, text width=2.4cm, align=center]
     at (-1.35, -3.58) {substituted span\\[-2pt]\emph{1--2 vs 6 tokens,\\[-2pt]never compared}};
\draw[decorate, decoration={brace, amplitude=3pt, mirror}, black!55]
     (0.03, -3.40) -- (1.72, -3.40);
\node[font=\scriptsize, anchor=north, text width=2.4cm, align=center]
     at (1.35, -3.58) {shared suffix\\[-2pt]\emph{verbatim; where\\[-2pt]they must agree}};

% ---------------------------------------------------------------- panel (b) ---
\begin{scope}[xshift=5.95cm, yshift=-0.95cm]
\node[font=\footnotesize\bfseries, anchor=west] at (-2.05, 1.77) {(b)};

\coordinate (hT)  at (-1.00,  0.80);
\coordinate (hTc) at (-1.40, -0.42);
\coordinate (hI)  at ( 1.10,  0.66);
\coordinate (hIc) at ( 1.50, -0.56);

\draw[thick, dash pattern=on 2pt off 1.6pt, black!60] (hT) -- (hTc)
      node[midway, left, font=\scriptsize, inner sep=1.5pt] {$d_{\mathrm{T}}$};
\draw[thick, dash pattern=on 2pt off 1.6pt, black!60] (hI) -- (hIc)
      node[midway, right, font=\scriptsize, inner sep=1.5pt] {$d_{\mathrm{I}}$};
\draw[line width=1.1pt, red!75] (hT) -- (hI)
      node[midway, above, font=\scriptsize, inner sep=2pt] {$d_{\mathrm{cross}}$};

\fill[orange!85] (hT)  circle (2.1pt);
\fill[orange!45] (hTc) circle (2.1pt);
\fill[blue!75]   (hI)  circle (2.1pt);
\fill[blue!40]   (hIc) circle (2.1pt);

\node[font=\scriptsize, anchor=east] at ($(hT)+(-0.10,0.02)$)  {$h(V_\mathrm{T})$};
\node[font=\scriptsize, anchor=east] at ($(hTc)+(-0.10,0)$)    {$h(V_\mathrm{T}')$};
\node[font=\scriptsize, anchor=west] at ($(hI)+(0.10,0.02)$)   {$h(V_\mathrm{I})$};
\node[font=\scriptsize, anchor=west] at ($(hIc)+(0.10,0)$)     {$h(V_\mathrm{I}')$};

\node[font=\footnotesize, anchor=north] at (0.15, -1.05)
     {$\displaystyle \mathrm{MSG}=
       \frac{d_{\mathrm{cross}}}{\tfrac{1}{2}(d_{\mathrm{T}}+d_{\mathrm{I}})}$};
\node[font=\scriptsize, anchor=north, text width=4.4cm, align=center]
     at (0.15, -2.00)
     {states at $m$: solid is the numerator,\\[-2pt]dashed the denominator};
\end{scope}

\end{tikzpicture}
\caption{\textbf{Cross-modal token substitutability, made well-posed.}
\textbf{(a)} Four views of one item share a prefix and a terminal suffix
verbatim and differ only in a substituted span: the span written out
($V_\mathrm{T}$), a same-modality re-expression of it ($V_\mathrm{T}'$), the
span rendered as glyphs ($V_\mathrm{I}$), and the same word in a second
typeface ($V_\mathrm{I}'$). The span occupies one to two text tokens but six
visual tokens (small boxes), so no positional correspondence exists along it
and we assert nothing about it. Equivalence is instead asserted at the
\emph{merge position} $m$, the first token of the shared suffix, located as
$\mathrm{len}(\text{sequence})-\mathrm{len}(\text{suffix tokens})$ and hence
model-agnostic --- it never needs to know how the processor expanded the image.
The strips are the actual stimulus at the frozen render geometry.
\textbf{(b)} The Modality Substitution Gap normalises the cross-modal distance
at $m$ by the within-modality control distances measured on the same items.
The denominator is what gives the ratio a reference value: it asks whether
crossing modalities costs more than re-expressing the same content inside one.
$\mathrm{MSG}\approx1$ means it does not.}
\label{fig:method}
\end{figure}
